%!TEX root = ./../../thesis-trier.tex

The notions of distinguishability of pairs of subsets and of complete distinguishability, and some accompanying results,
were introduced in~\cite{DBLP:conf/dlt/Hoffmann23}. 
Bondar \& Volkov~\cite{BondarV16,BondarCV23} introduced 
completely reachable automata.
\cref{prelim:lem:compl_reach_implies_SCC_sync,sync_max:cor:single-non-permutation}
were stated by them, 
see also~\cite[Sect.~6]{BondarCV23},~\cite{DBLP:conf/latin/CasasV22}
or~\cite[Prop.~3]{HoffmannCircularExtended}
for other statements.
A preliminary version of \cref{sync_max:cor:orbits-excluded-states-defect-one}
has been given in~\cite[Prop.~3]{DBLP:conf/lata/Hoffmann21a}.
\Cref{sync_max:sec:BVC-graphs} will present additional
results on completely reachable automata by Bondar \& Volkov~\cite{BondarV16,BondarV18}
and Bondar et al.~\cite{BondarCV23}.
The remaining results on reachable subsets from \cref{sync_max:sec:reachable} 
have not been mentioned in the literature before. 
In~\cref{sync_max:ex:k-reachable-not-sync}, we gave non-synchronizing $n$-state automata where every
$k$-set is reachable with $n = 2k - 2$. We do not know whether there exist such automata
where we have $n = k^2 - k - 1$ and $k \ge 2$, i.e., whether 
the bound from~\cref{sync_max:prop:reachability-k-sets-sync}
is optimal.
Complete reachability has been used in~\cite{HoffmannLATAextended2} to characterize primitive permutation groups (these results are also presented in \cref{sync_max:perm:sec:aut_sync_words} here).

%Maslennikova~\cite{DBLP:journals/corr/Maslennikova14,
%DBLP:journals/ijfcs/Maslennikova19} introduced the reset complexity
%of a given regular ideal language 
%as the size of a minimal automaton having it as its set of synchronizing words.
%Thus, \cref{sync_max:thm:lower_bound_SynA} 
%gives a lower bound for the reset complexity 

The distinguishability relation was introduced in~\cite{DBLP:conf/dlt/Hoffmann23,Hoffmann2024}.
In~\cite[Def.~5]{DBLP:journals/jalc/FrigeriR25} the first definition equals the indistinguishably relation
derived from our notion of distinguishability, but this has not been noted in~\cite{DBLP:journals/jalc/FrigeriR25}.

	\paragraph{An Erroneous Result.} In~\cite[Thm.~7]{DBLP:conf/dcfs/Hoffmann21} 
	% später sagen (and used in~\cite{dcfs12} ) 
	it was wrongly claimed that a transformation monoid is completely distinguishable
	if and only if this is true for all pairs
	of distinct $2$-sets (a related false statement
	appears in \cite[Lem.~3.1]{DBLP:conf/lata/Hoffmann21}). 
	\Cref{sync_max:sec:furtherexamples} lists various counterexamples.
	Now, the main result
	of~\cite[Thm.~12]{DBLP:conf/dcfs/Hoffmann21} (which is still true) characterizing primitive
	permutation groups in terms of the state complexity of associated
	automata used this false statement in its proof,
	as well as a special case of it stated before in~\cite[Prop.~4.8]{DBLP:conf/lata/Hoffmann21}.
	In \cref{sync_max:perm:sec:aut_sync_words}, \cref{sync_max:thm:prim},
    we present a correct proof, another correct proof has also been given in the extended version~\cite[Thm.~6.3]{HoffmannLATAextended2} 
	of the conference contributions~\cite{DBLP:conf/lata/Hoffmann21,DBLP:conf/dcfs/Hoffmann21}.	
	The false claim has also been used in~\cite{DBLP:conf/dcfs/Hoffmann22},
	and one main result from this paper,
	stating that a completely reachable SI-automaton
	is completely distinguishable,
	is in fact false, 
	see \cref{sync_max:fig:SCC_not_sufficient} for a counterexample.
	In the paper~\cite{DBLP:conf/dcfs/Hoffmann22}, it was shown 
	that in a completely reachable SI-automaton
	all pairs of distinct $2$-sets are distinguishable,
	and then the wrong claim has been used to conclude
	complete distinguishability.	
	The inductive ``proof'' in~\cite[Thm.~7]{DBLP:conf/lata/Hoffmann21a} (not contained in the proceedings version)
	missed a case (we note that this is related
	to the condition $\applytrans{\delta}{A}{u} \ne \applytrans{\delta}{B}{u}$
	in \crefitem{sync_max:thm:char_distinguishable}{dist:item:dcps}), which
	was reported to us by an anonymous reviewer of a paper that was (of course)
	rejected. The counterexample in \cref{sync_max:fig:2-sets-but-not-all:left} is due to the reviewer.
	We genuinely thank the reviewer for spotting the mistake!
	% We are really thankful, 
	%as we would probably still believe in this wrong argument.
	Our attempts at fixing it led to more insight,
	the notion of strong $\frac{3}{2}$-connectivity, and involvement of the theory of vertex-transitive graphs.